\section{Discussion}

\subsection{What actually drives the choice}

The descriptor that organises the winner map is $\NumMechFeature = p\,D /
\mathrm{cap}_C$: the share of the shortest corridor's capacity that the guided
cohort alone would consume. It is a product of demand and penetration, and
neither alone reproduces the structure --- which is why the best single-factor
rule reaches only \NumWinMapOne{} against \NumWinMapConst{} for a constant
rule. Demand determines whether the corridor is under pressure; penetration
determines how much of that pressure a policy commands. A system routing 20\% of
traffic in a saturated network and one routing 80\% in a lightly loaded one are
different decision problems, and the descriptor says so.

That is the answer to RQ2: the preference boundary is explicable in mechanistic
terms, and the explanation is physical rather than statistical. It is also why
the boundary must be quoted as a bracket. The grid samples
$\NumMechFeature$ at discrete levels, and the transition sits between
\NumMechBracketLo{} and \NumMechBracketHi{}. A point estimate would be a
property of the fitting procedure, not of the traffic.

\subsection{Does machine learning add value over the mechanistic rule?}

Yes, and the reason is instructive. The mechanistic rule is \emph{more accurate
and more expensive} than doing nothing: it names the best policy more often than
the fixed policy does (\NumBOneTopOne{} against \NumBZeroTopOne{}) and costs
more (\NumBOneRegret~s against \NumBZeroRegret~s). The learned selector is
\emph{less accurate} than the logistic baseline and cheaper than everything
except the retrospective optimum.

The difference is not model capacity. It is what the model is asked to predict.
B1, B2 and B3 predict the identity of the winner and are trained on a
classification loss that treats every context alike. B4 predicts the
\emph{advantage} of each policy over the fixed default, so a context where the
margin is 20\,s contributes a hundred times the gradient of one where it is
0.2\,s. Under an \NumAsymRatio$\times$ asymmetry between the fixed policy's
upside and its downside, that weighting is the whole game. This is the
predict-then-optimise argument~\cite{elmachtoub2022spo}, and it appears here not
as a refinement but as the difference between a selector that helps and one that
hurts.

The practical corollary for anyone building such a system: report decision
regret, and do not report selection accuracy as though it were a proxy for it.
In this study the two rank the methods differently.

\subsection{When adaptation is useful, and when a fixed policy is preferable}

Adaptation helps in-distribution and hurts outside it. On the interpolation
control the learned selector cuts regret from \NumOODThreeMidBZero~s to
\NumOODThreeMidBFour~s; on unseen high demand it raises it from
\NumOODOneHighBZero~s to \NumOODOneHighBFour~s, and under an unseen disruption
regime from \NumOODSevenIncBZero~s to \NumOODSevenIncBFour~s. The value of
adaptation is a function of regime coverage, not of model capacity, and no
amount of capacity substitutes for having seen the regime.

But the decisive observation is one of magnitude. The best deployable selector
saves \NumBFourVsBZero~s per vehicle against the best fixed policy, on a mean
journey time of \NumFixedCostPThree~s. Choosing the wrong fixed policy costs
\NumFixedWorstPenaltyPct{}. The two decisions differ by nearly three orders of
magnitude, and only one of them is worth an operator's attention.

\subsection{What the abstention mechanism is really telling us}

B5 abstains in \NumAbstainRate{} of contexts and never improves on the fixed
policy. It would be easy to present this as a failure of the gate, or to loosen
it until it acted. It is neither. The conformal interval on the predicted
advantage is wider than the advantage, so a gate that requires a positive lower
bound cannot certify a gain --- and should not. The gate is doing exactly what a
calibrated instrument is supposed to do when the effect is smaller than the
uncertainty.

The cost of that correctness is the in-distribution gain: B5 forgoes the
\NumBFourVsBZero~s that B4 achieves. That is the honest trade, and it is the
answer to RQ4. A selective mechanism prevented every out-of-distribution harm we
constructed, and it did so by declining to act at all.

\subsection{Implications for traffic operators}

Three, in decreasing order of value.

First, \textbf{choose the fixed policy deliberately}. In this environment the
gap between the best and worst fixed policy is \NumFixedWorstPenaltyPct{} of
system journey time. That is where the engineering effort belongs.

Second, \textbf{decide which criterion is being optimised, because the policies
specialise}. The load-balancing policy minimises journey time in
\NumSpecCOneN{} of \NumContextsAnalysed{} contexts; the reliability-aware policy
minimises stopped delay in \NumSpecCThreeN{}. An operator whose objective is
queue clearance at signalised junctions and one whose objective is corridor
journey time should not deploy the same policy, and this is not a preference
weighting --- it is a different argmin in \NumBestCThreeDiffersPct{} of
contexts.

Third, \textbf{do not build a per-context selector for this decision}. It works,
it is reproducible, and it is worth less than a tenth of a per cent. If a
selector is built anyway, give it a support gate and an abstention path, because
its out-of-distribution behaviour is worse than doing nothing.

\subsection{Why this is a useful negative result}

The literature on route guidance contains a persistent expectation that
context-adaptive selection should pay. What this study shows is that whether it
pays is an empirical property of the portfolio, not a general one, and that it
can be measured before any selector is built --- by comparing the best fixed
policy with the retrospective best, and comparing that gap with the noise floor
of the same comparison. Both quantities are cheap. Neither requires a model.

Had that comparison been run first here, it would have shown a
\NumHeadroomPct{} gap sitting below the resolution of a five-seed design, and
the selector would not have been built. We suggest it as a standing first step:
\emph{measure the headroom and the noise before choosing a model class.}
